\documentclass[12pt]{article}
\usepackage[utf8]{inputenc}
\usepackage[T1]{fontenc}
\usepackage{hyperref}
\title{Yanda Language Resources for K-12 Curriculum Development}
\author{Compiled by Claude Sonnet 5.5 (claude-sonnet-5-5) for Andr\'as Kornai}
\date{}
\begin{document}
\maketitle

\section{Introduction}
Yanda (ISO 639-3 yda) is an extinct Australian language of Queensland. It is probably Pama-Nyungan, possibly Maric, and may have been close to Guwa. It is known only from a single 1886 wordlist of about 110 words, and the transcription of that list is reported to be highly inaccurate. Search results showed no speakers, no teaching materials and no revitalization programme. The similarly named Yandruwandha (a different, better-documented language) must not be confused with it. The searches did not establish whether the two names refer to related groups.

\section{Primary Educational Resources}
None were found. No primers, textbooks or school programmes for Yanda appeared in the search results.

\section{Linguistic and Reference Materials}
\begin{itemize}
\item Wikipedia, ``Yanda language'': a summary of the single wordlist and the uncertain classification. \url{https://en.wikipedia.org/wiki/Yanda_language}
\item Endangered Languages Project, bibliography entry (lang/6741), which appeared in the Yanda search. I did not verify which language this entry covers. \url{https://endangeredlanguages.com/lang/6741/bibliography?hl=fr}
\item For disambiguation only, Wikipedia on the distinct language Yandruwandha, with Gavan Breen's 2004 grammar as its main documentation. \url{https://en.wikipedia.org/wiki/Yandruwandha_language}
\end{itemize}

\section{Religious and Community Resources}
None were found.

\section{Digital and Audio Resources}
No audio, video or software was found.

\section{Curriculum Development Context}
There is essentially no usable material. A K-12 heritage curriculum could at most be built around the 1886 wordlist, with clear caveats about its unreliable transcription. It should also draw on the documentation of neighbouring languages, and should be developed with the local Aboriginal communities concerned.
\end{document}